\documentclass[12pt]{article}
\usepackage[utf8]{inputenc}
\usepackage[T1]{fontenc}
\usepackage{hyperref}
\title{Southwestern Tamang Language Resources for K-12 Curriculum Development}
\author{Compiled by Claude Sonnet 5 (claude-sonnet-5) for Andr\'as Kornai}
\date{}
\begin{document}
\maketitle

\section{Introduction}
Southwestern Tamang (ISO 639-3 tsf) is spoken mainly in the southern plains (terai) of Nepal, especially around Chitwan. Very little was found for this language. Most Tamang resources online are for Western Tamang (tdg) or other Tamang varieties, which are separate languages. They are listed below only as clearly labelled related material, not as Southwestern Tamang resources.

\section{Primary Educational Resources}
None found. No Southwestern Tamang primer, textbook or mother-tongue education programme appeared in the searches.

\section{Linguistic and Reference Materials}
No dictionary or grammar for Southwestern Tamang was found. For comparison only, there is a Western Tamang--Nepali dictionary (a different language, tdg) made with SIL International and Mother Tongue Center Nepal: \url{https://www.webonary.org/wtamang/overview/?lang=en}.

\section{Religious and Community Resources}
An Asia Harvest people-group profile says no Scripture, Jesus film or recordings exist for Southwestern Tamang. The profile is undated and may be out of date: \url{https://www.people-groups.asiaharvest.org/India/TamangSouthwestern.pdf}. For Western Tamang (a different language), Scripture is listed at \url{https://www.scriptureearth.org/00eng.php?iso=tdg} and recordings at \url{https://globalrecordings.net/en/language/tdg}.

\section{Digital and Audio Resources}
None found for Southwestern Tamang itself.

\section{Curriculum Development Context}
The searches turned up no existing school materials for this variety. Any curriculum would need new documentation work, perhaps using the Western Tamang dictionary only as a model. Checking Ethnologue, Scripture Earth and Nepali mother-tongue education authorities directly is recommended.
\end{document}